\documentclass{article}
\usepackage{amsmath,amssymb}
\usepackage{xurl}
\usepackage[hidelinks]{hyperref}
% Shared commands for notes in docs/paper-gaps/.

\DeclareUnicodeCharacter{2080}{\ifmmode{}_{0}\else\textsubscript{0}\fi}
\DeclareUnicodeCharacter{2081}{\ifmmode{}_{1}\else\textsubscript{1}\fi}
\DeclareUnicodeCharacter{2082}{\ifmmode{}_{2}\else\textsubscript{2}\fi}
\DeclareUnicodeCharacter{2099}{\ifmmode{}_{n}\else\textsubscript{n}\fi}

\newcommand{\Lean}{\textsc{Lean}}
\ExplSyntaxOn
\NewDocumentCommand{\leanid}{m}{
  \ifmmode
    \textsf{\detokenize{#1}}
  \else
    \group_begin:
    \urlstyle{sf}
    \tl_set:Nn \l_tmpa_tl {#1}
    \tl_replace_all:Nnn \l_tmpa_tl {\_} {_}
    \exp_args:NV \path \l_tmpa_tl
    \group_end:
  \fi
}
\ExplSyntaxOff
\providecommand{\tr}{\operatorname{tr}}
\newcommand{\tnleanIssueBase}{https://github.com/LionSR/TNLean/issues/}
\newcommand{\tnleanPrBase}{https://github.com/LionSR/TNLean/pull/}
\newcommand{\tnleanDocBase}{https://sirui-lu.com/TNLean/docs/}
\newcommand{\ghissue}[1]{\href{\tnleanIssueBase#1}{GitHub issue \##1}}
\newcommand{\ghpr}[1]{\href{\tnleanPrBase#1}{PR \##1}}
\newcommand{\doclink}[1]{\href{\tnleanDocBase#1.html}{\path{#1}}}

% Dirac notation and the identity operator, as in blueprint/src/macros/common.tex.
% \providecommand keeps note-local definitions authoritative.
\usepackage{dsfont}
\providecommand{\ket}[1]{|#1\rangle}
\providecommand{\bra}[1]{\langle #1|}
\providecommand{\braket}[2]{\langle #1 | #2 \rangle}
\providecommand{\ketbra}[2]{|#1\rangle\!\langle #2|}
\providecommand{\Id}{\mathds{1}}
\providecommand{\one}{\mathds{1}}
\providecommand{\C}{\mathbb{C}}
\providecommand{\MN}[1]{M_{#1}(\C)}
\providecommand{\MD}{M_D(\C)}

% Verdict marker: \gapnote{<kind>}{<status>}. Kind is the policy
% classification (clarification, local-correction, scope-restriction,
% unfaithful, false-source, open-gap); status is open, wip (elimination
% underway), resolved, or historical. Machine-read by the site index;
% prints a verdict line.
\newcommand{\gapnote}[2]{\par\noindent\textbf{Verdict:} #1 (#2).\par}

\title{The Tree Circuit as a Multiscale Network: Two-Site Injectivity and Chain Length}
\author{}
\date{2026-09-30}
\begin{document}
\maketitle
\gapnote{scope-restriction}{open}

\section{What the source states}
Let $A$ be a normal tensor of physical dimension $d$ and bond dimension $D$, and
let $\ket{\phi_N}$ be the normalized translation-invariant state it generates on
a ring of $N$ sites. The paragraph ``Connection to MERA''
of~\cite{Malz2024Preparation} reads the tree circuit of eq.~(16) as a
finite-range multiscale entanglement renormalization ansatz with
$O(\log\log N)$ layers: the isometries $V^{(i)}$ of eq.~(16) are the isometries
of the network, every disentangler is the identity except those of the first
layer, which prepare the fixed-point state $\ket\Omega$, and hence, within the
approximation error $\varepsilon$, normal translation-invariant matrix product
states are such networks with $O(\log\log N)$ layers.

The paragraph ``The tree-RG circuit'' of the source describes the layers as
isometries from dimension $D^2$ to $D^4$, except the lowest, and its picture
starts from a blocked tensor. A layer $V^{(i)}$ is the partial isometry of the
polar decomposition of a two-site blocked map; it is an isometry exactly when
that map is injective. For $d^2<D^2$ the lowest layer, a map from $\mathbb{C}^{D^2}$ to
$(\mathbb{C}^d)^{\otimes2}$, cannot be an isometry at all.

\section{The restriction adopted}
The formal statements\footnote{Formal declarations:
    \leanid{MPSPreparation.approximatingMPVState_eq_state_treeMERA},
    \leanid{MPSPreparation.exists_state_treeMERA_approximationError_le}, and
    \leanid{MPSPreparation.exists_state_treeMERA_approximationError_le_of_lt}
    in \doclink{TNLean/MPS/Preparation/TreeMERA}.}
assume two things.
\begin{enumerate}
    \item The two-site blocked tensor of $A$ is injective. Then every layer
          $V^{(i)}$ is an isometry: blocking the two-site blocked tensor
          $V^{(1)}T_1$ twice gives $(V^{(1)})^{\otimes2}$ applied to $T_1$
          blocked twice, so the two-site blocked tensor of $T_1$ is injective,
          and so on up the tree.
    \item The chain length is $N=M2^{k+1}$, so that the $M$ blocks of the
          approximating state of eq.~(10) all have $q=2^{k+1}$ sites and the
          tree of each block is complete.
\end{enumerate}
Under these assumptions the approximating state of eq.~(10) is exactly the state
of the network with $k+1$ isometry layers, and its error against $\ket{\phi_N}$
is at most $\varepsilon$ once $2^{k+1}\ge\xi\log(CN/\varepsilon)$, with $\xi$
the correlation length of $A$ and $C$ depending only on $A$. The least such $k$
satisfies $k<\log_2(\xi\log(CN/\varepsilon))$, which is the source's
$O(\log\log(N/\varepsilon))$ count.

\section{Elimination}
\begin{enumerate}
    \item \emph{Two-site injectivity.} A normal tensor has an injectivity length
          $L_0$ independent of $N$. The tensor $A$ blocked over $L_0$ sites has
          an injective one-site, hence two-site, blocked tensor, is again normal,
          and generates $\ket{\phi_{NL_0}}$ from $\ket{\phi_N}$ of the blocked
          tensor after the regrouping of sites. Applying the statements to it
          gives networks on sites of dimension $d^{L_0}$, a constant, which is
          what the source's picture, starting from a blocked tensor, shows.
          What remains is to transport the gauge data (left-canonical form, the
          fixed point $\sigma$ and the bound on the subleading eigenvalues,
          raised to the power $L_0$) to the blocked tensor.
    \item \emph{Chain length.} The unequal blocks of the depth bound for every
          chain length (the note on divisible lengths) let the last block have
          length $q'$ with $q\le q'<2q$. A complete binary tree has $2^{k+1}$
          leaves, so the last block needs either one extra isometry layer
          $\mathbb{C}^{D^2}\to(\mathbb{C}^d)^{\otimes r}$ with $r\le 3$ at its lowest level, or a
          network whose last site is coarser; either extends the definition of
          the network and is left open.
\end{enumerate}

\bibliographystyle{alpha}
\bibliography{references}
\end{document}
